\lecture{8}{Thu 12 Feb 2026 14:01}{Lol brian taught class tday}

To compute the cohomology $H^i_{\Lie }(\mathfrak{g} ,M)$, we need to compute $\operatorname{Ext}^n_{U\mathfrak{g} } (\mathbb{K} ,M)$. Since $\mathbb{K} $ never changes, it might be easiest to find a free resolution of $\mathbb{K} $ and then to tensor with $M$ over $U\mathfrak{g} $ and take the homology of that. We proceed by doing just this using the Chevalley-Eilenberg resolution. To do this, we work with dg Lie algebras.

\section{dg Lie $\mathbb{K} $-algebras}

Recall that a graded $\mathbb{K} $-vector space is a collection $\left\{ V^i \right\}_{i\in\mathbb{Z} } $ of $\mathbb{K} $-vector spaces. It is differential graded if it has a differential $d^i : V^i \to V^{i+1} $ for all $i$ with $d^2=0$. The graded tensor product of graded $\mathbb{K} $-modules is
\[
	M\otimes_\mathbb{K} N \coloneqq \bigoplus_{i+j=k} M^i\otimes _\mathbb{K} N^j
.\]

\begin{definition}[dg Lie $\mathbb{K} $-algebra]\label{3.3.1}
	A differential graded Lie $\mathbb{K} $-algebra is a differential graded $\mathbb{K} $-vector space $\mathfrak{g} $ together with a bracket $[-,-] : \mathfrak{g} \otimes_\mathbb{K} \mathfrak{g}  \to \mathfrak{g} $ such that
	\begin{itemize}
		\item The bracket respects degrees:
			\[
				[\mathfrak{g}_i, \mathfrak{g} _j]\subseteq \mathfrak{g} _{i+j} ;
			\]
		\item the bracket satisfies \emph{graded anticommutivity}:
			\[
				[x,y]=-(-1)^{\left| x \right| \left| y \right| } [y,x]
			;\]
		\item the differential is a derivation for the bracket:
			\[
				d[x,y]=[dx,y]+(-1)^{\left| x \right| } [x,dy]
			;\]
		\item the bracket satisfies the graded Jacobi identity:
			\[
				[x,[y,z]]=[[x,y],z]+(-1)^{\left| x \right| \left| y \right| } [y,[x,z]
			.\]
	\end{itemize}
	Here if $x\in \mathfrak{g} $ is a homogeneous element, then $\left| x \right| $ is the \emph{degree} of $x$, meaning the number $i$ such that $x\in \mathfrak{g} _i$.
\end{definition}

\begin{example}\label{3.2.2}
	Any ordinary Lie $\mathbb{K} $-algebra can be viewed as a dg Lie $\mathbb{K} $-algebra concentrated in degree 0. More interestingly, given a Lie $\mathbb{K} $-algebra $\mathfrak{g} $, define the \emph{de Rham} dg Lie $\mathbb{K} $-algebra $\mathfrak{g} _{\mathrm{dR} } $ which is concentrated in degrees 0 and $-1$, and has $\mathfrak{g} =\mathfrak{g} _{\mathrm{dR} }^0=\mathfrak{g}_{\mathrm{dR} }^{-1}$. Let $x,y\in \mathfrak{g} ^0$ and $x',y'\in \mathfrak{g} ^1$. Then the bracket is defined as
	\[
		[x+x',y+y']_{\mathfrak{g} _{\mathrm{dR} } }=[x,y]_{\mathfrak{g} } +[x,y']_{\mathfrak{g} } +[x',y]_{\mathfrak{g} } 
	.\]
	The differential is given by the identity $1 : \mathfrak{g}_{\mathrm{dR} } ^{-1} \cong \mathfrak{g} _{\mathrm{dR} } ^0$.
\end{example}

\begin{remark}\label{3.2.3}
	There is a quasi-isomorphism $\mathfrak{g} _{\mathrm{dR} } \simeq 0$.
\end{remark}

\begin{definition}\label{3.2.4}
	A \emph{differential graded commutative $\mathbb{K} $-algebra} $A$ is a differential graded $\mathbb{K} $-module which is
	\begin{itemize}
		\item a graded $\mathbb{K} $-algebra, meaning it has a multiplication map $C\otimes_\mathbb{K} C\to C$ which satisfies $A^iA^j \subseteq A^{i+j} $ and graded commutivity:
			\[
				xy=(-1)^{\left| x \right| \left| y \right| } yx
			;\]
		\item a differential graded associative $\mathbb{K}$-algebra, meaning that the differential acts as a derivation:
			\[
				d(xy) = d(x)y + (-1)^{\left| x \right| } xd(y)
			.\]
	\end{itemize}
\end{definition}

\begin{remark}\label{3.3.5}
	Let $A$ be a dg commutative $\mathbb{K} $-algebra. Then $A$ is a dg Lie $\mathbb{K} $-algebra by taking the Lie bracket to be a sort of graded commutator. Define
	\[
		[a,b]\coloneqq ab-(-1)^{\left| a \right| \left| b \right| } ba
	.\]
	We show this satisfies the axioms. Graded antisymmetry:
	\[
		[b,a]=ba-(-1)^{\left| a \right| \left| b \right| } ab=-(-1)^{\left|  \right| \left| b \right| } \left( ab-(-1)^{\left| a \right| \left| b \right| } ba \right) =(-1)^{\left| a \right| \left| b \right| } [a,b]
	.\]
	The graded Jacobi identity is more than we would like to type, so we omit it. Compatibility with the differential:
	\[
		d[a,b]=d(ab) - (-1)^{\left| a \right| \left| b \right| } d(ba) = d(a)b + (-1)^{\left| a \right| } ad(b) - (-1)^{\left| a \right| \left| b \right| } \left( d(b)a + (-1)^{\left| b \right| } bd(a) \right) 
	\]
	\[
		=\left( d(a)b - (-1)^{(\left| a \right| +1)\left| b \right| } bd(a) \right) +(-1)^{\left| a \right| } \left( ad(b)-(-1)^{\left| a \right| (\left| b \right| +1)} d(b)a \right)
	\]
	\[
		=\left( d(a)b-(-1)^{\left| da \right| \left| b \right| } b(da) \right) +(-1)^{\left| a \right| } \left( ad(b)-(-1)^{\left| a \right| \left| db \right| } d(b)a \right) =[da,b]+(-1)^{\left| a \right| } [a,db]
	.\]
	Here we used that $\left| a \right| =\left| a \right| ^2$. Given any dg commutative algebra $A$, we denote the corresponding dg Lie $\mathbb{K} $-algebra by $\Lie (A)$. Note that $\Lie $ is a functor $\mathcal{D} \mathcal{G} \Alg_{\mathbb{K} } \to \mathcal{D} \mathcal{G} \Lie _\mathbb{K} $.
\end{remark}

\begin{definition}\label{3.3.6}
	Let $\mathfrak{g} $ be a dg Lie $\mathbb{K} $-algebra and $A$ a dg commutative $\mathbb{K} $-algebra. Then the \emph{universal enveloping algebra} $U\mathfrak{g} $ is universal in the sense that there are isomorphisms
	\[
		\mathcal{D} \mathcal{G}\mathrm{Lie }_\mathbb{K} (\mathfrak{g} ,\Lie (A)) \cong \mathcal{D} \mathcal{G} \Alg_{\mathbb{K} } (U\mathfrak{g} ,A)
	\]
	natural in $\mathfrak{g} $ and $A$. By construction, $U\dashv \Lie $.
\end{definition}

\begin{definition}\label{3.3.7}
	Let $M$ be a dg $\mathbb{K} $-module. Then the \emph{graded symmetric $\mathbb{K} $-algebra} is the free graded commutative $\mathbb{K} $-algebra on $M$, and is given by formal symmetric products of elements in $M$. Explicitly, if $T(M)$ is the tensor algebra,
	\[
		\operatorname{Sym} (M) \coloneqq \frac{T(M)}{\left\langle x\otimes y - y\otimes x \right\rangle } 
	.\]
\end{definition}

\begin{remark}\label{3.3.8}
	The symmetric algebra of a dg $\mathbb{K} $-module is canonically a dg commutative $\mathbb{K} $-algebra with the following differential. First, the tensor algebra $T(M)$ on a dg $\mathbb{K} $-module $(M,d)$ is a dg associative $\mathbb{K} $-algebra with differential
	\[
		d(v_1 \otimes \cdots \otimes v_n) = \sum_{i=1}^{n} (-1)^{\left| v_1 \right| +\cdots +\left| v_{i-1}  \right| } v_1 \otimes \cdots \otimes dv_i \otimes \cdots \otimes v_n
	.\]
	An annoying check shows that the differential indeed preserves the quotient, meaning we obtain a differential on $\operatorname{Sym} (M)$ given by
	\[
		d(v_1 \cdots v_n) = \sum_{i=1}^{n} (-1)^{\left| v_1 \right| + \cdots + \left| v_{i-1}  \right| } v_1 \cdots dv_i \cdots v_n
	.\]
\end{remark}



\begin{theorem}[PBW]\label{3.3.9}
	Let $\mathfrak{g} $ be a dg Lie algebra. Then $U\mathfrak{g}\cong \operatorname{Sym} (\mathfrak{g} )$. 
\end{theorem}

\begin{lemma}\label{lma:asjkskjas}
	$\mathcal{U} (\mathfrak{g} _{\mathrm{dR} } )$ is a free $\mathcal{U} (\mathfrak{g} )$-module.
\end{lemma}
\begin{proof}
	The universal enveloping algebra is monoidal $(\Mod _\mathfrak{g} ,\oplus )\to (\Mod _{\mathcal{U} (\mathfrak{g} )} ,\otimes_\mathbb{C}  )$. Additionally, $\mathcal{U} (\mathfrak{g} \oplus \mathfrak{h} )\cong \mathcal{U} (\mathfrak{g} )\otimes_\mathbb{C} \mathcal{U} (\mathfrak{h} )$ is free (take a basis for $\mathfrak{h} $). Apply this to $\mathfrak{g} _{\mathrm{dR} } $ to obtain via PBW that
	\[
		\mathcal{U} (\mathfrak{g} _{\mathrm{dR} } )\cong \mathcal{U} (\mathfrak{g} )\otimes _\mathbb{C} \operatorname{Sym} (\mathfrak{g} [1])
	\]
	as $\mathcal{U} (\mathfrak{g} )$-modules. This proves it.
\end{proof}

Since $\mathcal{U} (\mathfrak{g} _{\mathrm{dR} }) \simeq \mathbb{C} $, we have a free resolution as a dg Lie algebra. Hence
\[
	\operatorname{Ext}_{\mathcal{U} (\mathfrak{g} )} ^i(\mathbb{C} ,M) \cong H^i\left( \Mod_{\mathcal{U} (\mathfrak{g} )} (\mathcal{U} (\mathfrak{g} _{\mathrm{dR} } ),M \right) 
.\]
If $A$ is a free $\mathbb{C} $-algebra, then $\Mod _A(A\otimes F,M)\cong \Mod _\mathbb{C} (F,M)$. Applying this here yields
\[
	\Mod _{\mathcal{U} (\mathfrak{g} )} (\mathcal{U} (\mathfrak{g} _{\mathrm{dR} } ),M)\cong \Mod_{\mathcal{U} (\mathfrak{g} )} (\mathcal{U}(\mathfrak{g} )\otimes_\mathbb{C} \operatorname{Sym} (\mathfrak{g} [1]),M) \cong \Vect_{\mathbb{C} } (\operatorname{Sym} (\mathfrak{g} [1]),M)
.\]
We often call this complex $C^{\bullet} (\mathfrak{g} ;M)$. In degree 0, then this is $\Vect_{\mathbb{C} } (\mathbb{C} ,M)\cong M$. In degree 1, we take the degree $-1$ generators $\Vect_{\mathbb{C} } (\mathfrak{g} ,M)$. In degree 2, its $\Vect_{\mathbb{C} } (\Lambda ^2\mathfrak{g} ,M)$. This generalizes.
\[
	C^i(\mathfrak{g} ;M)=\Vect_{\mathbb{C} } (\Lambda ^i\mathfrak{g} ,M)
.\]

Cheveleakslky-Eilenberg's theorem si that if $G$ is a connected compact Lie group, then the cohomology of the invariants of the de Rham complex $\Omega (G)^G$ are that of $C^{\bullet}(\mathfrak{g} ;\mathbb{C} ) $. More powerfully, $\Omega (G)^G \cong C^{\bullet} (\mathfrak{g} ;\mathbb{C} )$. This follows because your tangent bundle for Lie groups is trivial so ur forms are constant coefficient. Hence we obtain a differential on $C^{\bullet} (\mathfrak{g} ;\mathbb{C} )$ via the exterior derivative.

We claim for any $M$, then $C^{\bullet} (\mathfrak{g} ;M)$ is a $C^{\bullet} (\mathfrak{g};\mathbb{C} )$ module. Take $C^{\bullet} (\mathfrak{g} ;M) \cong C^{\bullet} (\mathfrak{g} )\otimes M$. We may now define the differential. For trivial coefficients, $d$ being a derivation implies we need only define it on generators:
\[
	C^{\bullet} (\mathfrak{g} ) \cong \operatorname{Sym} (\mathfrak{g} [1]^*) = \Lambda ^{\bullet} \mathfrak{g} ^*
\]
has it as $0 : \mathbb{C}  \to \mathfrak{g} $, and it suffices now only to define it on $\mathfrak{g}^* \to \Lambda^2 \mathfrak{g}^* $ since all others are generated on it. We define this as the adjunct of $[-,-]: \Lambda ^2\mathfrak{g}  \to \mathfrak{g} $. Check wikipedia for a full explicit definition.

For $C^{\bullet} (\mathfrak{g} ;M)$, we only need it on degree 0 since it is a $C^{\bullet} (\mathfrak{g} )$-module. The map $M\to \Hom(\mathfrak{g} , M) $ is given by $m\mapsto (x\mapsto xm)$. Wikipedia also computes the first few cohomology groups. H2 corresponds to ``central extensions'' modulo some sort of isomorphism classes. These are extensions of $\mathfrak{g} $ as a Lie algebra by $\mathbb{C} $, meaning Lie algebras $\mathfrak{h} $ which extend $\mathfrak{g} $ in the sense taht
\[\begin{tikzcd}[row sep = 2.5em, column sep = 2.5em]
0\ar[r]  & \mathbb{C} \ar[r]  & \mathfrak{h} \ar[r]  & \mathfrak{g} \ar[r]  & 0
\end{tikzcd}\]
is exact. Short exact sequences of vector spaces always split, so $\mathfrak{h} \cong \mathfrak{g} \oplus \mathbb{C} $. Thus its bracket has to be $[x,y]+\varphi (x,y)c$ for some basis $c\in\mathbb{C} $ and some $\varphi : \Lambda ^2\mathfrak{g}  \to \mathbb{C} $. What in all of the fucks is happening????? review these notes later.  see wikipedia page on Lie group cohomology, nLab page on dg Lie algebras, \cite{cg1} A.3, and check the paper they cite at the end and apparently Lurie treats this for $\infty $-categories somewhere so maybe I can find that.
